% !TEX program = xelatex
\documentclass[UTF8,fontset=fandol,a4paper,11pt]{ctexart}
\usepackage{../../_manual_common/hysim_manual}
\title{\textbf{HySim-XJTU-HRPES}\\[0.4em]{\Large server 模块运行时契约手册}\\
{\large 会话修订、ETag、异步作业与版本路由}}
\author{西安交通大学 HRPES 团队}
\date{\today}
\begin{document}
\maketitle
\begin{abstract}
本文档对应 \srcpath{src/server/runtime\_api\_v1.cpp}、\srcpath{runtime\_api\_v1.hpp} 与
\srcpath{edition\_profile.cpp}。本模块是 HTTP/JSON 编排层，不拥有新的电力系统数学模型；本文形式化的
只有代码真实执行的修订一致性、缓存、作业和路由分类规则。
\end{abstract}
\tableofcontents

\section{模块定位、输入输出与边界}
\texttt{RuntimeApiV1} 接收案例构建器和工作线程数，向 \texttt{httplib::Server} 注册多会话接口。
输入是 HTTP 请求、JSON 和当前会话；输出是带状态码、schema、revision 与有效性字段的 JSON。
GUI 后端 \srcpath{tests/run\_gui\_server.cpp} 不属于本目录，二者不能共享会话或缓存。

\section{会话修订和 ETag 的代码等价模型}
每个 \texttt{ApiSession} 保存系统指针和无符号修订号 $r$，其默认值由
\srcpath{src/server/runtime\_api\_v1.cpp:ApiSession} 初始化为 0。ETag 由
\srcpath{model\_etag}（同文件:model\_etag）精确构造为
\begin{equation}
E(s,r)=\texttt{"hysim-}+s+\texttt{-r}+\operatorname{to\_string}(r)+\texttt{"}.
\end{equation}
创建会话时若请求体包含模型，代码把 $r$ 置为 $1$（:register\_routes）；PUT 成功替换模型后执行
$r\leftarrow r+1$（:register\_routes）。读请求若 \texttt{If-None-Match}=E(s,r) 返回 304
（:check\_if\_none\_match）。写请求的前置条件由 \srcpath{require\_if\_match}（:require\_if\_match）实现：
\begin{equation}
\texttt{If-Match存在}\land \texttt{If-Match}=E(s,r).
\end{equation}
缺少头返回 428；值不相等返回 412；两种情况均不得修改系统或修订号。

\section{拓扑缓存一致性}
会话缓存保存 $r_c$ 与 LOD2 拓扑 DOM。\srcpath{session\_topology\_lod2}
（\srcpath{src/server/runtime\_api\_v1.cpp:ApiSession}）的更新条件为
\begin{equation}
\texttt{cache为空}\lor r_c\ne r.
\end{equation}
条件为真时重新序列化并令 $r_c\leftarrow r$，否则按引用复用缓存。LOD0/1 从 LOD2 聚合；viewport 必须
同时提供 $x_{min},y_{min},x_{max},y_{max}$，缺一即抛出无效参数。子图查询的
\fld{max\_nodes} 由 \srcpath{query\_int} 钳制在 $[1,5000]$，默认 500
（:topology\_subgraph），BFS 达到该数量即停止扩张（:topology\_subgraph）。

\section{异步作业与陈旧结果}
\texttt{ApiJob} 保存提交时的 $r_j$ 与系统快照（\srcpath{src/server/runtime\_api\_v1.cpp:ApiJob}）。
\srcpath{job\_json}（:job\_json）返回结果时执行
\begin{equation}
stale=(r_{current}\text{存在})\land(r_{current}\ne r_j).
\end{equation}
因此旧作业可以完成，但必须在响应中携带 \fld{model\_revision}、
\fld{current\_model\_revision} 与陈旧标志；绝不能把旧结果标成当前模型结果。代码中的作业状态字符串和
取消行为是权威状态机，HTTP 202 表示已接受而不是已完成，删除成功返回 204。

\section{版本路由分类}
\texttt{EditionRouteAccess} 的三个实际枚举值是 \texttt{Retained}、\texttt{Disabled}、
\texttt{Unclassified}（\srcpath{src/server/edition\_profile.hpp:EditionRouteAccess}）。
\srcpath{classify\_trial\_route}（\srcpath{src/server/edition\_profile.cpp:classify\_trial\_route}）依次处理 OPTIONS、
非 API 静态内容、两组禁用路由、保留路由与 v1 模式；其余 API 返回
\texttt{Unclassified}。\texttt{Disabled} 必须拒绝，\texttt{Unclassified} 不能自动当作
\texttt{Retained}。版本 profile 不是授权系统，也不改变底层求解器能力。

\section{算法、复杂度、异常与并发}
会话/作业查找平均为哈希表 $O(1)$；拓扑全量序列化为 $O(N+E)$，子图 BFS 为 $O(N+E)$ 且受
\fld{max\_nodes} 截断。互斥量保护会话修订和缓存；耗时分析使用提交时快照，避免持锁求解。不存在会话、
模型为空、revision 冲突、参数缺失、作业不存在、后端不支持和求解失败分别使用代码定义的 HTTP 错误，
不应统一伪装成 200。

\section{接口关系、验证与工业指标}
本模块依赖 model、io、api 和各分析模块；HTTP schema 不应反向渗透进求解核心。测试需覆盖 304、428、
412、并发写冲突、缓存随 revision 失效、作业陈旧标记、取消竞争、LOD/viewport 边界和 disabled route。
工业指标包括错误状态码准确率、陈旧结果漏报率（目标 0）、revision 单调性、并发数据竞争数（目标 0）、
拓扑分页内存峰值与 P95 延迟。

\section{适用范围与局限}
本手册不描述 GUI 单会话路由，不承诺分布式持久化、跨进程作业恢复或身份认证。当前缓存仅属进程内会话；
进程重启后的状态不在模型内。后续改进应由 OpenAPI/schema 测试自动生成字段矩阵，并对状态转换做属性测试。
\section{跨模块工业级技术评价基线}

本章规定本手册所述模块进入工程算例、批量研究或生产服务前必须共同满足的评价基线。
具体模型公式、变量、约束和专用算法以正文相应章节为准；本章不扩大源码已经实现的范围。

\subsection{输入、输出、边界条件与数据结构审计}
输入审计应逐项检查有限性、单位、上下界、枚举值和引用完整性。系统对象必须先按所需层级完成
校验；AC 与 DC 母线采用域限定映射，组件稳定 \texttt{index}、向量位置和临时图索引不得混用。
输出审计必须记录单位、索引空间、模型范围、收敛或可行状态、后端、容差、迭代/节点计数和告警。
空系统、空时域、孤岛、重复 ID、零容量、退化约束与不可用可选后端均属于边界测试集。

\subsection{工程假设与数值稳定性分析}
模型使用的平衡、线性化、稳态、独立性、概率分布、时间离散或网络等值假设必须与应用场景匹配。
数值评价至少包含变量缩放、绝对/相对残差、可行性违反、条件数或枢轴质量、步长/阻尼、迭代停滞
和随机种子复现性。若采用正则化、平滑、回退、松弛或近似，应同时报告触发条件、参数值、对主指标
的敏感性以及是否改变模型口径；不得用“已返回结果”替代收敛或有效性证明。

\subsection{复杂度与资源分析}
令 $n_x$、$n_c$、$n_t$、$n_s$、$|V|$、$|E|$ 分别表示变量、约束、时段、场景、节点和边数。
报告应分别给出建模、求解、后处理的时间与内存。图遍历通常为 $O(|V|+|E|)$；逐时段/逐场景
流程至少线性增长为 $O(n_t n_s)$ 次子问题调用；稠密线性代数最坏可达 $O(n_x^3)$，稀疏方法则应
报告结构非零数、填充率和因子分解峰值内存；MILP/NLP 不给出虚假的多项式最坏界，而报告节点数、
间隙、时限和实测扩展曲线。复杂度结论必须基于固定硬件、固定后端和可复现算例族。

\subsection{异常工况处理}
非法输入应在进入求解器前返回结构化错误；不可行、无界、不收敛、时限、内存不足和后端缺失必须
相互区分。部分结果只有在对应 \fld{ValidityFlags}、\fld{model\_scope}、
\fld{model\_limitations} 或等价字段允许时才可使用。异常路径不得静默丢弃 AC/DC 资产、制造平衡
节点、把 NaN 改写为零或把近似解标为全局最优；系统原始对象应保持可恢复和可重试。

\subsection{与其他模块的接口关系}
接口链路遵循“工程语义模型 $\to$ 校验 $\to$ 投影/装配 $\to$ 求解或分析 $\to$ 结果回投影”。
跨模块传递时保留稳定 ID、单位、符号、时间基准和场景身份；消费潮流、OPF、动态或可靠性结果前，
先检查其收敛、覆盖范围和有效性标志。HTTP/JSON/Python 层只做语义保持适配，不重新解释求解结果。

\subsection{测试与验证建议}
最低验证矩阵包括：手算小例、标准算例、边界/异常输入、随机性质测试、算法或后端交叉校核、
序列化往返、稳定 ID 回投影、AC/DC 同号母线、重复运行确定性和规模扩展测试。数值算法还应加入
残差独立重算、雅可比有限差分或自动微分校核；近似模型应与高保真模型比较并给出误差分布，而非
只报告单个平均值。回归报告必须列明构建配置、算例版本、容差、通过/失败/跳过数及未验证范围。

\subsection{工业级评估指标}
\begin{itemize}
  \item \textbf{正确性}：守恒残差、约束最大违反、解析/基准误差、结果回投影一致性；
  \item \textbf{鲁棒性}：收敛率、不可行识别率、异常分类准确率、扰动敏感性与随机复现率；
  \item \textbf{性能}：端到端时延、吞吐量、峰值内存、并行效率与规模增长斜率；
  \item \textbf{可追溯性}：输入模式版本、稳定 ID 覆盖率、模型范围/限制字段完整率、告警可定位率；
  \item \textbf{工程有效性}：与标准、外部引擎或现场数据的偏差，以及误判/漏判对业务指标的影响。
\end{itemize}
指标应预先给出验收阈值和失败处置；未达到阈值时只能标记为实验性或受限适用。

\subsection{适用范围、局限性与后续改进}
本手册只评价当前源码覆盖的模型、后端和数据结构，不对未建模物理过程、未安装求解器或未执行的
外部验证作保证。后续改进应按“缺口 $\to$ 可量化指标 $\to$ 源码/测试变更 $\to$ 独立验证”闭环，
优先消除会造成错误结论的语义缺口，其次提升数值鲁棒性与性能，最后扩展模型范围。

\end{document}
